\documentclass[10pt,a4paper]{amsart}
\usepackage[margin=19mm]{geometry}
\usepackage{amsmath,amssymb,mathtools}
\usepackage[scr=rsfs,cal=euler]{mathalfa}
\usepackage{xfrac}
\usepackage[unicode,hidelinks,bookmarks=false]{hyperref}
\newcommand{\ie}{i.e.\ }
\newcommand{\cf}{cf.\ }
\newcommand{\resp}{respectively\ }
\newcommand{\Subsection}{\subsection}
\providecommand{\nobreakdash}{\nobreak\hbox{-}}
\setlength{\emergencystretch}{2em}
\multlinegap0pt
\usepackage{bm}
\usepackage{bbm}
\usepackage{dsfont}
\usepackage{xspace}
\usepackage{cancel}
\usepackage{mathtools}
\usepackage{cleveref}
\usepackage{amsthm}
\makeatletter
\@ifundefined{theorem}{\newtheorem{theorem}{Theorem}}{}
\makeatother
\newcommand{\Batrix}[1]{\bm{\Matrix{#1}}}
\newcommand{\Hilb}{\mathrm{H}}
\newcommand{\Index}[1]{{#1}}
\newcommand{\Matrix}[1]{\mathsf{#1}}
\newcommand{\Norm}[3][]{#1\| #2 #1\|_{\mathrm{#3}}}
\newcommand{\N}{\mathbb{N}}
\newcommand{\pinv}{\dagger}
\newcommand{\trans}{\intercal}
\newcommand{\truncate}[3][]{#1\lfloor #2 #1\rfloor_{#3}}
\title[Matrices over a Hilbert space and their low-rank cross appro]{Matrices over a Hilbert space and their low-rank cross approximation}
\author{Stanislav Budzinskiy}
\date{Source: arXiv:2505.05134v2 (CC BY 4.0)}
\begin{document}
\maketitle

\noindent\emph{Source and licence.} Matrices over a Hilbert space and their low-rank cross approximation. Version arXiv:2505.05134v2 (CC BY 4.0), distributed under the Creative Commons Attribution 4.0 International licence, as recorded in the arXiv licence metadata for this exact version and shown on the abstract page https://arxiv.org/abs/2505.05134v2. Full rights notice: licenses/arxiv-2505.05134-CC-BY-4.0.txt.\par\smallskip

\noindent\textit{Editorial scope.} Counted unit: the Theorem whose source label is \texttt{theorem:cross\_error\_maxvol} in the article (source0.tex lines 1014--1027, source label \texttt{theorem:cross\_error\_maxvol}), delivered together with the article's own complete original proof of it (source0.tex lines 1028--1076, one \texttt{proof} environment of 49 lines). No mathematics is written, repaired or re-derived: statement and proof are the article's own text line for line. Required context retained uncounted: theorem:cgr\_error\_maxvol (lines 801--810). Every definition, display and label of those lines is retained unchanged. Omitted from this package and not redistributed: 1--800, 811--1013, 1077--1270. Nothing inside a retained excerpt is omitted or altered: every retained body is delivered exactly as the article's lines are written, so no omission receipt of any kind accompanies this package. Citations: the retained excerpts carry no citation command at all, so no bibliography entry of the article is transcribed and no reference is redistributed. Reference adaptation: none was needed and none was made; every reference inside the delivered unit resolves inside it, so no unresolved reference or citation is produced. Printed slips kept literal: no printed word, symbol, hypothesis or number of the counted statement or of its proof was altered. Layout and class: the article's own class is \texttt{siamart250211} with packages including amsfonts, amssymb, amsopn, eucal, bm, upgreek, graphicx, subcaption, epstopdf, booktabs, algorithm2e; the wrapper is the lane's pinned \texttt{amsart} editorial wrapper at 10pt on A4, which restates the theorem-like environments the retained text needs and adds the article's own macro definitions that the retained text needs (\texttt{\textbackslash Batrix}, \texttt{\textbackslash Hilb}, \texttt{\textbackslash Index}, \texttt{\textbackslash Matrix}, \texttt{\textbackslash N}, \texttt{\textbackslash Norm}, \texttt{\textbackslash pinv}, \texttt{\textbackslash trans}, \texttt{\textbackslash truncate}), with extra wrapper packages (bm, bbm, dsfont, xspace, cancel, mathtools, cleveref). The delivered numbering is the unit's own. Assets: no retained span contains a figure environment, a graphics-inclusion command, a PDF-page inclusion, a table or a TikZ picture; no image file is copied into this package, the package declares no asset dependency and no image file is redistributed with it. Licence: version arXiv:2505.05134v2 of this article is distributed under the Creative Commons Attribution 4.0 International licence, as recorded in the arXiv licence metadata for this exact version and shown on the abstract page https://arxiv.org/abs/2505.05134v2. A full rights notice is delivered at licenses/arxiv-2505.05134-CC-BY-4.0.txt. Characters: every retained excerpt is pure ASCII, so the delivered engine, XeLaTeX with its default Latin Modern text font, reports no missing character.
\begin{theorem}
\label{theorem:cgr_error_maxvol}
Let $\Batrix{A} \in \Hilb^{m \times n}$ and $\kappa \in \N$.
Let $\Batrix{G}$ have the maximum $\kappa$-projective volume among all $|\Index{I}| \times |\Index{J}|$ submatrices of $\Batrix{A}$.
If $\nu_{\kappa}(\Batrix{G}) > 0$,
\begin{equation*}
    \Norm{\Batrix{A} - \Batrix{C} \truncate{\Batrix{G}}{\kappa}^\pinv \Batrix{R}}{\ell_{\infty}(\Hilb)} \leq \sqrt{1 + \frac{\kappa - \Delta_\kappa'}{|\Index{I}| + 1 - \kappa + \Delta_\kappa'}} \sqrt{1 + \frac{\kappa}{|\Index{J}| + 1 - \kappa}} \sigma_{\kappa + 1}(\Batrix{A})
\end{equation*}
with $\Delta_\kappa' = \min_{i \not\in \Index{I}} \min_{1 \leq j \leq n} \Delta_{\kappa}\big( \Batrix{A}(\Index{I} \cup \{i\}, \Index{J} \cup \{j\}) \big)$.
\end{theorem}

\begin{theorem}
\label{theorem:cross_error_maxvol}
Let $\Batrix{A} \in \Hilb^{m \times n}$. 
Let $\rho,\kappa \in \N$ and $\theta_c, \theta_r \geq 1$.
Suppose $\Batrix{G}$ has the $\theta_c$-quasi-maximum $\kappa$-projective volume among all $|\Index{I}| \times |\Index{J}|$ submatrices of $\Batrix{A}$ and $\Batrix{G}^\trans$ has the $\theta_r$-quasi-maximum $\rho$-projective volume among all $|\Index{J}| \times |\Index{I}|$ submatrices of $\Batrix{C}^\trans$.
Denote $\Batrix{X} = (\truncate{\Batrix{G}^\trans}{\rho}^\pinv \Batrix{C}^\trans)^\trans \Batrix{G} (\truncate{\Batrix{G}}{\kappa}^\pinv \Batrix{R})$.
If $\nu_{\kappa}(\Batrix{G}) > 0$ and $\nu_{\rho}(\Batrix{G}^\trans) > 0$,
\begin{align*}
    \Norm{\Batrix{A} - \Batrix{X}}{\ell_{\infty}(\Hilb)} &\leq \theta_c \sqrt{1 + \frac{\kappa - \Delta_\kappa'}{|\Index{I}| + 1 - \kappa + \Delta_\kappa'}} \sqrt{1 + \frac{\kappa}{|\Index{J}| + 1 - \kappa}} \sigma_{\kappa + 1}(\Batrix{A}) \\
    %
    &+ \theta_r \theta_c \sqrt{1 + \frac{\rho}{|\Index{I}| + 1 - \rho}} \sqrt{1 - \frac{1}{\theta_c^2} + \frac{\kappa}{|\Index{J}| + 1 - \kappa}} \sigma_{\rho + 1}(\Batrix{C}^\trans).
\end{align*}
with $\Delta_\kappa' = \min_{i \not\in \Index{I}} \min_{1 \leq j \leq n} \Delta_{\kappa}\big( \Batrix{A}(\Index{I} \cup \{i\}, \Index{J} \cup \{j\}) \big)$.
\end{theorem}

\begin{proof}
By triangle inequality, we have
\begin{equation*}
    \Norm{\Batrix{A} - \Batrix{X}}{\ell_{\infty}(\Hilb)} \leq \Norm{\Batrix{A} - \Batrix{C} \truncate{\Batrix{G}}{\kappa}^\pinv \Batrix{R}}{\ell_{\infty}(\Hilb)} + \Norm{\Batrix{C} \truncate{\Batrix{G}}{\kappa}^\pinv \Batrix{R} - \Batrix{X}}{\ell_{\infty}(\Hilb)}.
\end{equation*}
We bound the first term via a quasi-maximum-volume variant of \cref{theorem:cgr_error_maxvol}, the proof of which is a direct modification of the original proof:
\begin{equation*}
    \Norm{\Batrix{A} - \Batrix{C} \truncate{\Batrix{G}}{\kappa}^\pinv \Batrix{R}}{\ell_{\infty}(\Hilb)} \leq \theta_c \sqrt{1 + \frac{\kappa - \Delta_\kappa'}{|\Index{I}| + 1 - \kappa + \Delta_\kappa'}} \sqrt{1 + \frac{\kappa}{|\Index{J}| + 1 - \kappa}} \sigma_{\kappa + 1}(\Batrix{A}).
\end{equation*}
Suppose that the maximum error in the second term is attained at the $(i,j)$th entry and denote $\Batrix{c}^\trans = \Batrix{C}(i,:)$ and $\Batrix{r} = \Batrix{R}(:, j)$ as previously:
\begin{align*}
    \Norm{\Batrix{C} \truncate{\Batrix{G}}{\kappa}^\pinv \Batrix{R} - \Batrix{X}}{\ell_{\infty}(\Hilb)} &= \Norm{(\Batrix{C} - (\truncate{\Batrix{G}^\trans}{\rho}^\pinv \Batrix{C}^\trans)^\trans \Batrix{G}) \truncate{\Batrix{G}}{\kappa}^\pinv \Batrix{R}}{\ell_{\infty}(\Hilb)} \\
    %
    &= \Norm{(\Batrix{c} - \Batrix{G}^\trans \truncate{\Batrix{G}^\trans}{\rho}^\pinv \Batrix{c})^\trans \truncate{\Batrix{G}}{\kappa}^\pinv \Batrix{r}}{\Hilb} \\
    %
    &\leq \Norm{\Batrix{c} - \Batrix{G}^\trans \truncate{\Batrix{G}^\trans}{\rho}^\pinv \Batrix{c}}{\ell_2(\Hilb)} \Norm{\truncate{\Batrix{G}}{\kappa}^\pinv \Batrix{r}}{\ell_2}.
\end{align*}
As in the proof of \cref{theorem:cgr_error_maxvol}, we consider different cases based on the position of $(i,j)$ relative to $(\Index{I}, \Index{J})$.
When $i \in \Index{I}$ and $\Batrix{c} = \Batrix{G}^\trans \Matrix{e}_{\omega}$,
\begin{align*}
    \Norm{\Batrix{c} - \Batrix{G}^\trans \truncate{\Batrix{G}^\trans}{\rho}^\pinv \Batrix{c}}{\ell_2(\Hilb)} &= \Norm{(\Batrix{G}^\trans - \Batrix{G}^\trans \truncate{\Batrix{G}^\trans}{\rho}^\pinv \Batrix{G}^\trans) \Matrix{e}_{\omega}}{\ell_2(\Hilb)} \\
    %
    &\leq \Norm{\Batrix{G}^\trans - \truncate{\Batrix{G}^\trans}{\rho}}{\ell_2 \to \ell_2(\Hilb)} = \sigma_{\rho + 1}(\Batrix{G}^\trans) \leq \sigma_{\rho + 1}(\Batrix{C}^\trans).
\end{align*}
When $i \not\in \Index{I}$, we repeat the argument from the third part of the proof of \cref{theorem:cgr_error_maxvol}.
Namely, we form an augmentation $\Batrix{M} = [\Batrix{G}^\trans~\Batrix{c}]$, take an auxiliary matrix $\Matrix{T}$, multiply $\Batrix{K} = \Batrix{M} \Matrix{T}$, obtain an equality
\begin{equation*}
    \nu(\Batrix{K}) = \nu_{\rho}(\Batrix{G}^{\trans}) \Norm{\Batrix{c} - \Batrix{G}^\trans \truncate{\Batrix{G}^\trans}{\rho}^\pinv \Batrix{c}}{\ell_2(\Hilb)},
\end{equation*}
and derive the bound
\begin{equation*}
   \Norm{\Batrix{c} - \Batrix{G}^\trans \truncate{\Batrix{G}^\trans}{\rho}^\pinv \Batrix{c}}{\ell_2(\Hilb)} = \frac{\nu(\Batrix{K})}{\nu_{\rho}(\Batrix{G}^{\trans})} \leq \frac{\nu_{\rho + 1}(\Batrix{M})}{\nu_{\rho}(\Batrix{G}^{\trans})} \leq  \theta_r \sqrt{1 + \frac{\rho}{|\Index{I}| + 1 - \rho}} \sigma_{\rho + 1}(\Batrix{C}^\trans).
\end{equation*}
When $j \in \Index{J}$ and $\Batrix{r} = \Batrix{G} \Matrix{e}_{\omega}$, we have
\begin{equation*}
    \Norm{\truncate{\Batrix{G}}{\kappa}^\pinv \Batrix{r}}{\ell_2} = \Norm{\truncate{\Batrix{G}}{\kappa}^\pinv \Batrix{G} \Matrix{e}_{\omega}}{\ell_2} \leq \Norm{\truncate{\Batrix{G}}{\kappa}^\pinv \Batrix{G}}{\ell_2 \to \ell_2} \leq 1.
\end{equation*}
When $j \not\in \Index{J}$, consider $\Batrix{M} = [\Batrix{G}~\Batrix{r}]$, and let $\truncate{\Batrix{G}}{\kappa} = \Batrix{U}_\kappa \Matrix{\Sigma}_\kappa \Matrix{V}_\kappa^\ast$ be an SVD.
Then, by the matrix determinant lemma,
\begin{align*}
    \nu_\kappa^2(\Batrix{M}) \geq \nu^2(\Batrix{U}_\kappa^\ast \Batrix{M}) &= \det\Big( \Matrix{\Sigma}_\kappa^2 + (\Batrix{U}_\kappa^\ast \Batrix{r}) (\Batrix{U}_\kappa^\ast \Batrix{r})^\ast \Big) \\
    %
    &= \nu_\kappa^2(\Batrix{G}) (1 + \Norm{\Matrix{\Sigma}_\kappa^{-1} \Batrix{U}_\kappa^\ast \Batrix{r}}{\ell_2}^2) = \nu_\kappa^2(\Batrix{G}) (1 + \Norm{\truncate{\Batrix{G}}{\kappa}^\pinv \Batrix{r}}{\ell_2}^2),
\end{align*}
and consequently
\begin{equation*}
    \Norm{\truncate{\Batrix{G}}{\kappa}^\pinv \Batrix{r}}{\ell_2} \leq \sqrt{\frac{\nu_\kappa^2(\Batrix{M})}{\nu_\kappa^2(\Batrix{G})} - 1} \leq \sqrt{\theta_c^2 + \frac{\theta_c^2 \kappa}{|\Index{J}| + 1 - \kappa} - 1} = \theta_c \sqrt{1 - \frac{1}{\theta_c^2} + \frac{\kappa}{|\Index{J}| + 1 - \kappa}}.
\end{equation*}
\end{proof}

\end{document}
